
\documentclass[11pt]{article}
\usepackage[utf8]{inputenc}
\usepackage{amsmath,amssymb,amsthm}
\usepackage{hyperref}
\usepackage{url}
\usepackage{cleveref}
\usepackage{geometry}
\usepackage{enumitem}
\usepackage{xcolor}

\geometry{margin=1in}

\newtheorem{theorem}{Theorem}[section]
\newtheorem{lemma}[theorem]{Lemma}
\newtheorem{proposition}[theorem]{Proposition}
\newtheorem{corollary}[theorem]{Corollary}
\theoremstyle{definition}
\newtheorem{definition}[theorem]{Definition}
\newtheorem{axiom}[theorem]{Axiom}
\theoremstyle{remark}
\newtheorem{remark}[theorem]{Remark}
\newtheorem{example}[theorem]{Example}

% Lean footnote helper: keep Lean references out of the main exposition.
% `\path{...}` allows line breaks in long file paths.
\newcommand{\LeanFile}[1]{\footnote{\textbf{Lean}: \path{#1}.}}

% Minimal-fix additions: both used below (line 62 and line 69) but were only
% defined inside the commented-out preamble block. Lifted into the active
% preamble so the file builds; commented-out block left untouched per "idk" note.
\newcommand{\R}{\ensuremath{\mathbb{R}}}
\newcommand{\draftnote}{\noindent\textit{Draft extended abstract for author revision.}}

% idk, I'll use what ks-math-foundations has above.
% \usepackage[margin=1in]{geometry}
% \usepackage[T1]{fontenc}
% \usepackage[utf8]{inputenc}
% \usepackage{amsmath,amssymb,amsthm}
% \usepackage{booktabs}
% \usepackage{array}
% \usepackage{enumitem}
% \usepackage{xcolor}
% \usepackage{hyperref}
% \usepackage{url}
% \hypersetup{colorlinks=true,linkcolor=blue!50!black,citecolor=blue!50!black,urlcolor=blue!50!black}
% \setlist[itemize]{leftmargin=1.4em,itemsep=0.25em,topsep=0.25em}
% \setlist[enumerate]{leftmargin=1.6em,itemsep=0.25em,topsep=0.25em}
% \newcommand{\WMPLN}{WM--PLN}
% \newcommand{\PLN}{PLN}
% \newcommand{\KS}{K\&S}
% \newcommand{\codepath}[1]{\path{#1}}
% \newcommand{\leanref}[1]{\path{#1}}
% \newcommand{\Evidence}{\ensuremath{\mathcal{E}}}
% \newcommand{\World}{\ensuremath{\mathcal{W}}}
% \newcommand{\Query}{\ensuremath{\mathcal{Q}}}
% \newcommand{\R}{\ensuremath{\mathbb{R}}}
% \newcommand{\draftnote}{\noindent\textit{Draft extended abstract for author revision.}}
% \sloppy

% Draft by Author.  This is meant as a draft for JoeFest 2026.
% Comments are largely directed to the beloved AI assistants <3.
\title{Formal Foundations of Inference in Lean\\From Finite Lattices to Infinitary Probability}
\author{Zarathustra Amadeus Goertzel}
\date{\today}
\begin{document}
\maketitle
\draftnote

\begin{abstract}
Knuth and Skilling's \emph{Foundations of Inference} aims to derive basic probability theory, Bayes' rule, and the entropy formula from order and symmetry constraints over distributive lattices.
This aims to provide clearer intuition as to why the laws of probability take the form they do, continuing the tradition of Cox's axioms.
I set out to formalize their work with the help of Codex and Claude Code\footnote{Most of the work was done by Codex 5.2 and Claude Code Opus 4.5.  I will use the pronoun 'we' in this document to fairly credit the AI assistance without which this would not have been done.}.
Curiously, a missing \emph{separation axiom} was discovered, without which the additive representation does not hold.
This was discovered to be equivalent, in the bounded-below setting, to Alimov's \emph{no-anomalous-pairs} condition (Alimov~1950)~--~the classical hypothesis under which Hölder's representation theorem yields an order-embedding of an ordered semigroup into~$\R$. %% Thank Claude for the correction here!  I am not sure to what extent we need the "in the bounded-below setting" part.  Claude, you must be absolutely careful not to introduce extraneous details.
These axioms play the role of the hidden richness assumptions that Cox-style approaches often require.
Regularity assumptions were also needed to derive the Entropy formula.
Additionally, we find that adding three completeness axioms ($\sigma$-closure of events, sequential completeness of the value scale, and Scott continuity of the valuation) recovers Kolmogorov's $\sigma$-additivity as a theorem~--~providing a unified foundation for both discrete and continuous probability theory. %% Thank you for the correctness as to the axioms we added.  The abstract is where we should insert the marketing claim and we should not understate the results.  I did not intend to cover all measure theory.  The point is to recover the core of Kolmogov's measure-theoretic probablitiy theory, and I tihnk your versoin UNDERSELLS what we have.  That is very, very bad.
This work both presents the formalization of the foundations of inference and showcases how AI-assisted formalization can assist in axiomatic ablation studies over formal systems. %% I probably did not write a single line of Lean.  To call it AI-assisted theorem proving with substanital human steering is probably to over-weigh my own role.

% As to the suggestion, no, I will not add childish line counts  file counts.  Longer isn't bette.
% Unifying Kolmogorov's probability theory with the finite versions is a strong contribution.
% Saying "without which it does not hold" already implies we proved that.
% Yes, can you elaborate on what Halpbern said about Cox?  the 1999 paper?  We could mention that when we discuss continuity.
% Great, yes, I will mention the proof route diversity and ispmilicyt.
% We only have an abstract.  You dn't need to nag me about the draftnote.
% I find the ablation of axioms in appendix C etc with entropy confusing, so presenting that will take care :).
\end{abstract}

\section{Introduction}
Probability theory can be directly defined via Kolmogorov's measure theoretic axioms [cite]; however these do not explain \emph{why} this is the correct way to reason about Uncertainty.
Cox's theorem [cite] derives the rules of probability from desiderata for ``degrees of belief'', but requires continuity and differentiability assumptions.

%% [edit-1 reverted now that we're switching to pub26.bib: \cite{KnuthSkilling2012} -> \cite{ksfoi12}, matching the key in pub26.bib.]
Knuth and Skilling~\cite{ksfoi12} presented an algebraic foundation of probabilistic inference, which is about combining and evaluating information.
Under certain symmetries (associativity, distributivity, etc), the familiar calculi emerge.
Under weakened assumptions, credal sets, intervals, and weaker calculi for reasoning about uncertainty emerge.  %% Not sure how to word this.

This suggests a vision of future work in sketching out a \emph{probability hypercube} of domain assumption-grounded calculi for reasoning about uncertainty.
E.g., dropping commutativity can lead to forms of quantum probability in a natural manner.

\section{Events, Order, and Valuations}
%% [edit-2: "distributive of events" -> "distributive lattice of events" -- missing word "lattice".]
%% [edit-3: footnote: $human(x)$ etc. wrapped in \text{...}. Reason: bare names in math mode render as italic concatenated h*u*m*a*n*x. The \text{} keeps the prose font and renders correctly as the predicate name "human(x)".]
We begin with a distributive lattice of events over a set $E$ where order lines up with entailment: if a entails b, then a should not be assigned a larger plausibility than b\footnote{E.g., $\text{human}(x) \Rightarrow \text{animal}(x)$ and $P(\text{human}(x)) \leq P(\text{animal}(x))$.}.
%% [edit-4: typo "represenents" -> "represents".]
%% [edit-5: typo "comibnation" -> "combination".]
The top $\top$ and bottom $\bot$ events represent True and False.  Join ($a \vee b$) represents combination of events and meet ($a \wedge b$) represents intersection of events.

\begin{definition}
A \emph{distributive lattice} is a set $E$ equipped with $\vee,\wedge$ such that both operations are
associative, commutative, idempotent, and satisfy absorption, and such that $\wedge$ distributes over
$\vee$ (equivalently $\vee$ distributes over $\wedge$).
\end{definition}

%% Keeping below as formatting nudge
\[
  \Theta(x \oplus y) = \Theta(x) + \Theta(y).
\]


%% [edit-6: switched from inline \thebibliography{99} block to bibtex via pub26.bib.
%%   - \bibliographystyle{plain} matches the alphabetical-by-author rendering of your old inline block.
%%     Switch to alpha / amsplain / unsrt later if you prefer different label style.
%%   - The 6 inline entries that were NOT in pub26.bib (Halpern1999, Halpern2003,
%%     FaginHalpernMegiddo1990, Cox1946, Kolmogorov1933, Jaynes2003) are no longer rendered.
%%     None of them are cited from the body yet, so this is benign for now -- but you'll
%%     need to add them to pub26.bib when you fill in the [cite] placeholders on
%%     lines ~91-92 (Kolmogorov1933, Cox1946) and any future Halpern1999 mention.
%%   - Build sequence is now: pdflatex -> bibtex -> pdflatex -> pdflatex.]
\bibliographystyle{plain}
\bibliography{pub26}
\end{document}
